% Options for packages loaded elsewhere
\PassOptionsToPackage{unicode}{hyperref}
\PassOptionsToPackage{hyphens}{url}
%
\documentclass[
  10pt,
]{article}
\usepackage{amsmath,amssymb}
\usepackage{iftex}
\ifPDFTeX
  \usepackage[T1]{fontenc}
  \usepackage[utf8]{inputenc}
  \usepackage{textcomp} % provide euro and other symbols
\else % if luatex or xetex
  \usepackage{unicode-math} % this also loads fontspec
  \defaultfontfeatures{Scale=MatchLowercase}
  \defaultfontfeatures[\rmfamily]{Ligatures=TeX,Scale=1}
\fi
\usepackage{lmodern}
\ifPDFTeX\else
  % xetex/luatex font selection
  % Native editor uses its default main font.
  % Native editor uses its default sans font.
  % Native editor uses its default monospaced font.
\fi
% Use upquote if available, for straight quotes in verbatim environments
\IfFileExists{upquote.sty}{\usepackage{upquote}}{}
\IfFileExists{microtype.sty}{% use microtype if available
  \usepackage[]{microtype}
  \UseMicrotypeSet[protrusion]{basicmath} % disable protrusion for tt fonts
}{}
\makeatletter
\@ifundefined{KOMAClassName}{% if non-KOMA class
  \IfFileExists{parskip.sty}{%
    \usepackage{parskip}
  }{% else
    \setlength{\parindent}{0pt}
    \setlength{\parskip}{6pt plus 2pt minus 1pt}}
}{% if KOMA class
  \KOMAoptions{parskip=half}}
\makeatother
\usepackage{xcolor}
\usepackage[margin=0.8in]{geometry}
\usepackage{color}
\usepackage{fancyvrb}
\newcommand{\VerbBar}{|}
\newcommand{\VERB}{\Verb[commandchars=\\\{\}]}
\DefineVerbatimEnvironment{Highlighting}{Verbatim}{commandchars=\\\{\}}
% Add ',fontsize=\small' for more characters per line
\newenvironment{Shaded}{}{}
\newcommand{\AlertTok}[1]{\textcolor[rgb]{1.00,0.00,0.00}{\textbf{#1}}}
\newcommand{\AnnotationTok}[1]{\textcolor[rgb]{0.38,0.63,0.69}{\textbf{\textit{#1}}}}
\newcommand{\AttributeTok}[1]{\textcolor[rgb]{0.49,0.56,0.16}{#1}}
\newcommand{\BaseNTok}[1]{\textcolor[rgb]{0.25,0.63,0.44}{#1}}
\newcommand{\BuiltInTok}[1]{\textcolor[rgb]{0.00,0.50,0.00}{#1}}
\newcommand{\CharTok}[1]{\textcolor[rgb]{0.25,0.44,0.63}{#1}}
\newcommand{\CommentTok}[1]{\textcolor[rgb]{0.38,0.63,0.69}{\textit{#1}}}
\newcommand{\CommentVarTok}[1]{\textcolor[rgb]{0.38,0.63,0.69}{\textbf{\textit{#1}}}}
\newcommand{\ConstantTok}[1]{\textcolor[rgb]{0.53,0.00,0.00}{#1}}
\newcommand{\ControlFlowTok}[1]{\textcolor[rgb]{0.00,0.44,0.13}{\textbf{#1}}}
\newcommand{\DataTypeTok}[1]{\textcolor[rgb]{0.56,0.13,0.00}{#1}}
\newcommand{\DecValTok}[1]{\textcolor[rgb]{0.25,0.63,0.44}{#1}}
\newcommand{\DocumentationTok}[1]{\textcolor[rgb]{0.73,0.13,0.13}{\textit{#1}}}
\newcommand{\ErrorTok}[1]{\textcolor[rgb]{1.00,0.00,0.00}{\textbf{#1}}}
\newcommand{\ExtensionTok}[1]{#1}
\newcommand{\FloatTok}[1]{\textcolor[rgb]{0.25,0.63,0.44}{#1}}
\newcommand{\FunctionTok}[1]{\textcolor[rgb]{0.02,0.16,0.49}{#1}}
\newcommand{\ImportTok}[1]{\textcolor[rgb]{0.00,0.50,0.00}{\textbf{#1}}}
\newcommand{\InformationTok}[1]{\textcolor[rgb]{0.38,0.63,0.69}{\textbf{\textit{#1}}}}
\newcommand{\KeywordTok}[1]{\textcolor[rgb]{0.00,0.44,0.13}{\textbf{#1}}}
\newcommand{\NormalTok}[1]{#1}
\newcommand{\OperatorTok}[1]{\textcolor[rgb]{0.40,0.40,0.40}{#1}}
\newcommand{\OtherTok}[1]{\textcolor[rgb]{0.00,0.44,0.13}{#1}}
\newcommand{\PreprocessorTok}[1]{\textcolor[rgb]{0.74,0.48,0.00}{#1}}
\newcommand{\RegionMarkerTok}[1]{#1}
\newcommand{\SpecialCharTok}[1]{\textcolor[rgb]{0.25,0.44,0.63}{#1}}
\newcommand{\SpecialStringTok}[1]{\textcolor[rgb]{0.73,0.40,0.53}{#1}}
\newcommand{\StringTok}[1]{\textcolor[rgb]{0.25,0.44,0.63}{#1}}
\newcommand{\VariableTok}[1]{\textcolor[rgb]{0.10,0.09,0.49}{#1}}
\newcommand{\VerbatimStringTok}[1]{\textcolor[rgb]{0.25,0.44,0.63}{#1}}
\newcommand{\WarningTok}[1]{\textcolor[rgb]{0.38,0.63,0.69}{\textbf{\textit{#1}}}}
\setlength{\emergencystretch}{3em} % prevent overfull lines
\providecommand{\tightlist}{%
  \setlength{\itemsep}{0pt}\setlength{\parskip}{0pt}}
\setcounter{secnumdepth}{-\maxdimen} % remove section numbering
\usepackage{fvextra}
\fvset{breaklines=true,breakanywhere=true,fontsize=\scriptsize}
\setlength{\emergencystretch}{3em}
\ifLuaTeX
  \usepackage{selnolig}  % disable illegal ligatures
\fi
\usepackage{bookmark}
\IfFileExists{xurl.sty}{\usepackage{xurl}}{} % add URL line breaks if available
\urlstyle{same}
\hypersetup{
  hidelinks,
  pdfcreator={LaTeX via pandoc}}

\author{}
\date{}

\begin{document}

\section{Finite Galois transport of ordinary zero-Hessian
currents}\label{finite-galois-transport-of-ordinary-zero-hessian-currents}

Reviewed partial theorem, 4 October 2026. The original
atom P06.2 remains open. No catalogue description, earlier proof,
receipt, model verdict, or hash is a mathematical premise of this
argument.

\subsection{Result and scope}\label{result-and-scope}

Let K be a complete nontrivially valued non-Archimedean field, let E/K
be any finite Galois extension, and let X be a good, separated, strict,
boundaryless \(K\)-analytic space, pure of dimension \(n\ge 1\).
Retain the actual analytic structure of \(X\), including its nilpotents. Let
\(h:X\to\mathbb{R}\) be finite and continuous and let
\(\pi:X_E\to X\) be scalar projection. The assertion proved below is

\[
 d'd''[h]=0\quad\Longleftrightarrow\quad d'd''[h\circ\pi]=0.       \tag{T}
\]

The tests on each side are \textbf{all} ordinary smooth, compactly
supported (n-1,n-1)-forms. The respective canonical CLD calibrations are
used. No ramification, residue-characteristic, perfectness, normality,
reducedness, unit separation, Radon-Hessian, or invariant-form-descent
hypothesis is imposed. All open sets, neighborhoods, and supports below
use the ordinary Hausdorff topology. An affinoid neighborhood contains
an ordinary open neighborhood; this interior is not its absolute
analytic interior.

This is a scalar-current transport theorem, not recognition of a general
continuous zero-current function as classically pluriharmonic. It
establishes the finite-Galois current-transport gate and has passed
an independent Xean review. It does not establish two-sign regularization,
ascent through the completion of an infinite algebraic extension, purely
inseparable ascent, transcendental ascent, or full P06.2. None of those
consequences is asserted here. The original task and completion criteria
are preserved verbatim in original-task.json and original-problem.txt,
rather than rewritten as a weaker task.

Write \(D=d'd''\) and \(t_f=-\log|f|\) for an analytic unit.
With a harmless fixed duality sign suppressed, current vanishing means

\[
 B_h(\alpha):=\int_X h\,(D\alpha)_{\mathrm{Rad},K}=0
 \quad\text{for every }\alpha\in A_c^{n-1,n-1}(X).             \tag{1}
\]

Only differentiated smooth top forms in this proof define Radon
measures. No order-zero assertion is made about the Hessian current of h
itself.

\subsection{Exact imported analytic
foundations}\label{exact-imported-analytic-foundations}

These are foundations imported from primary texts, not conclusions from
earlier campaign manuscripts. The appendix embeds the primary bodies
used, including their complete endings and proofs when the primary text
gives one. The precise restricted uses are recorded here. The new
algebraic, Banach, support, and collar arguments are proved below.

\textbf{A. Affinoids and ordinary neighborhoods.} CLD,
arXiv:1204.6277v2, 25 July 2025, sections 11.1.1-2, 11.1.5-10 and
11.7.7: affinoid spectra are compact Hausdorff; an affinoid algebra is
noetherian with closed ideals and has surjective
spectrum-to-prime-spectrum map. Its completed residue field is the
completion of the fraction field of the quotient by the point's kernel.
Rational affinoid neighborhoods can be chosen inside any prescribed
ordinary open with every defining inequality strict at the point. Good
strict spaces over a nontrivially valued field have strict affinoid
neighborhood bases. Affinoid domains admit finite rational coverings
(11.1.6). Products of affinoids use completed tensor products, and
scalar projection is surjective. CLD 11.4.4 makes a finite algebra with
its finite-module Banach topology affinoid. Actual quotient Banach
norms, not spectral seminorms on reductions, are used below.

\textbf{B. Dimensions and boundaries.} CLD 11.2.2-3: separatedness
implies Hausdorffness; on a strict affinoid analytic dimension is Krull
dimension; analytic domains preserve local dimension; local dimension at
a point is the maximum dimension of the minimal-prime components
containing it. CLD 11.4.1 gives the composition equality for relative
interiors when the middle-to-target morphism is locally separated. In
its application below that morphism is the finite separated map
M(F)-\textgreater M(K). Berkovich, \emph{Etale cohomology for
non-Archimedean analytic spaces}, Publ. Math. IHES 78 (1993),
Proposition 1.5.5(i), (iii), (iv): an analytic domain's relative
interior is its ordinary topological interior. Under base change,
interior points remain interior; only this inclusion and the resulting
stability of boundarylessness are used, including for complete
ground-field extension.
Thus ordinary opens and finite scalar extensions of boundaryless spaces
are boundaryless. We prove componentwise purity separately.

\textbf{C. Ordinary forms.} CLD Definition 6.2.1 and 6.2.2 define
ordinary smooth forms by smooth tropical descriptions on ordinary
neighborhoods; these form ordinary sheaves stable under wedge, J, d',
d'\,'. CLD Proposition 15.1.1(b) gives descriptions using finitely many
analytic units. CLD 3.3.3 permits a coefficient form to be represented
by a Euclidean smooth form on an ambient open; this uses Lemma 1.9.4
only with the paralinéaire subset closed in that open. CLD 3.3.11-12
gives Leibniz and functorial pullback for forms. CLD Proposition 3.4.4
says that a coefficient-form pullback is zero exactly when its
restriction to the image of every compact domain is zero. These inputs
identify the ordinary sheaves on a same-carrier restriction of scalars,
without any claim about descent of arbitrary invariant forms.

\textbf{D. Compact cutoffs.} CLD 15.2.2 and Corollary 15.2.5: on a good
Hausdorff analytic space, for a compact C inside an ordinary open U
there is an ordinary smooth function equal to one near C with compact
support contained in U. We derive the finite compact partition and
continuous approximation statements below. We do not use the overly
broad printed compact-support assertion of CLD 15.2.9.

\textbf{E. Polyhedrality, calibration, and top measures.} Antoine
Ducros, \emph{Espaces de Berkovich, polytopes, squelettes et théorie des
modèles}, Confluentes Math. 4 (2012), 1250007, Theorem 3.2(1): for a
compact, Gamma-strict k-analytic space of dimension d and finitely many
analytic units, their multiplicative norm image is a c-polytope of
dimension \textbf{at most} d, with c=(Q,(\textbar k\^{}*\textbar{}
Gamma)\^{}Q). Taking minus logarithms gives the finite affine polyhedral
description used here. We apply it to strict affinoids (Gamma=\{1\}); no
equality of moment dimension and analytic dimension is used. The printed
universal equality in CLD 12.2.4 is not imported.

CLD 5.4.4 makes compact moment charts pleasant and identifies their
tropical polyhedron with the top-dimensional part of the image. CLD
14.3.2-5 defines the positive standard canonical calibration. On a
skeletal n-cell C, a torus projection q injective on its moment image
gives the constant calibration d\_C times the inverse image of the
primitive volume, where d\_C is the actual finite-flat moment degree.
CLD 14.5.1-2 pushes these calibrations to the moment image: on an n-cell
sigma its weight is the sum of the positive d\_C above sigma times the
corresponding primitive volume vector. CLD 5.8.1-5 associates a Radon
measure to every smooth top form, locally using precisely those
calibrations, with absolute variation equal to the associated Borel
measure. CLD 5.10.1(c),(d) gives the corresponding pushforward and
integration formula in a pleasant chart. Lower-dimensional cells carry
no n-dimensional top measure. All degrees count the actual finite fibre
algebras and their Artinian lengths. We spell out the resulting
variation estimate below.

\textbf{F. Two different scalar comparisons.} CLD Proposition 14.3.6 and
Proposition 15.3.1: for a complete extension L/K and a locally
integrable (n,n)-form omega,

\[
 \pi_* (\pi^*\omega)_{\mathrm{Rad},L}
       =\omega_{\mathrm{Rad},K}.                             \tag{2}
\]

There is \textbf{no {[}L:K{]} factor}. CLD Remark 14.4.3: for a finite
extension F/K and an F-analytic space viewed as a K-analytic space on
the \textbf{same carrier}, the dimensions and skeletal subsets agree and

\[
 \mu_K=[F:K]\mu_F.                                          \tag{3}
\]

These are different comparisons. Neither imports a pullback of {[}h{]}
or D{[}h{]}. The finite Artinian-algebra reason for the normalization in
(2) is reconstructed below, in addition to the embedded primary proof.

\textbf{G. Stokes.} CLD Theorem 15.7.1 gives a Hausdorff pure
n-dimensional analytic space its tropical structure with its analytic
boundary; 15.7.2 makes restrictions to analytic domains compatible. CLD
Proposition 5.11.1(a) gives Stokes for compact C1 (n-1,n)-forms. On a
boundaryless space their d' integral is zero, and the J-transformed
statement gives the d'\,' version. Applying these twice below proves the
integration by parts formula for arbitrary complementary forms. The
symmetric-form restriction in the separately printed Green formula is
not used.

The hypotheses just listed allow singularities and nilpotents. Their use
does not replace the measure by a normalized measure on the reduction.

\subsection{1. Actual finite tensors, fibres, and component
dimensions}\label{actual-finite-tensors-fibres-and-component-dimensions}

\subsubsection{1.0. Elementary field facts used in the
construction}\label{elementary-field-facts-used-in-the-construction}

For completeness, no primitive-element or fixed-field result is hidden
in the induction. A finite separable extension L/K of an infinite field
has {[}L:K{]} embeddings into an algebraic closure: along a finite
generating tower, each existing embedding extends in exactly as many
ways as the degree of the next separable minimal polynomial, by choosing
its distinct roots. If a1,\ldots,aq generate L, choose theta=sum\_j c\_j
a\_j so all these embeddings give distinct values of theta. This is
possible: each pair of embeddings forbids a nonzero linear equation on
(c\_j); putting (c\_j)=(1,t,\ldots,t\^{}\{q-1\}) makes each forbidden
equation a nonzero polynomial in t, so finitely many forbidden roots
cannot exhaust infinite K. The minimal polynomial of theta has at least
{[}L:K{]} distinct roots and degree at most {[}L:K{]}, hence
L=K{[}theta{]}. Our K is infinite because a finite field cannot have a
nontrivial absolute value.

For normal separable E/K, every embedding of an intermediate field
extends to an embedding of E by the same successive-root construction;
normality puts the image back in E. Thus all such embeddings extend to
G. A G-fixed element has only one K-conjugate, so is in K: E\^{}G=K. For
a subgroup H, put F=E\^{}H and choose a K-primitive b for E. Then
E=F{[}b{]}, and the product of T-hb over H has coefficients in F. Every
hb is a distinct root of the F-minimal polynomial of b, whereas that
minimal polynomial divides the product. Its degree is
\textbar H\textbar, proving E/F is Galois with group H. F/K is finite
separable as a subextension of E/K. For any embedding
tau:F-\textgreater E, its extension g in G identifies the stage E/tau(F)
with the conjugate subgroup gHg\^{}\{-1\}.

\subsubsection{1.1. Finite-module topology and
invariants}\label{finite-module-topology-and-invariants}

For a K-affinoid Banach algebra A and a finite extension E/K of degree
m,

\[
 A\widehat\otimes_K E=A\otimes_K E.                          \tag{4}
\]

Here is the topology justification. A fixed basis of a
finite-dimensional normed vector space over a complete valued field
identifies its norm with a norm equivalent to the coefficient maximum
norm. Induct on dimension: the span of all but the last basis vector is
complete, hence closed; the last vector has positive distance from it.
Scaling that distance bounds its coefficient in a vector; the
ultrametric inequality and induction bound the remaining coefficients.
The upper bound follows directly from the triangle inequality. Thus all
coordinate maps E-\textgreater K are bounded. Tensor them with A. The
tensor-norm upper bound in the other direction identifies the topology
on A tensor E with that on the finite direct sum A\^{}m. This direct sum
is complete, proving (4). This proof is for the actual algebra; no
quotient by its nilradical occurs.

For G=Gal(E/K), the action 1 tensor g on (4) is bounded. The linear
equations (g-1)v=0 on E have common kernel K. Tensoring this finite
kernel calculation with a K-vector space is exact; hence

\[
 (A\otimes_K E)^G=A.                                       \tag{5}
\]

No averaging in characteristic p is used. In particular, an invariant
unit, together with its invariant inverse, is a unit of A. For every
unit f of A tensor E,

\[
N(f)=\prod_{g\in G}g f\in A^\times.                         \tag{6}
\]

For clarity, the point action and the translates of unit domains used
later are analytic on the restricted K-carrier: (4) is finite over A, so
foundation A makes it a K-affinoid algebra, and 1 tensor g is a bounded
K-algebra automorphism. Its K-spectrum is the same as its E-spectrum
because every seminorm restricts to the unique valued-field norm on E.
The actual rational algebras and their unit germs agree. In particular
\textbar g f(y)\textbar=\textbar f(g\^{}\{-1\}y)\textbar. No invariance
of a continuous current or semilinear averaging of forms is used.

\subsubsection{1.2. Analytic fibres and their decomposition
groups}\label{analytic-fibres-and-their-decomposition-groups}

Let x belong to X, put C=H(x), and choose y above it. By (4) the actual
analytic fibre is M(C tensor\_K E). A primitive element b for the finite
separable extension E/K identifies this algebra with C{[}T{]}/P(T),
where P is separable. Its factorization is a finite product of fields;
each factor is a finite complete valued extension of C and has one
analytic point. The uniqueness of its extending absolute value can be
checked without residue assumptions: two such values are equivalent
norms on the finite C-vector space, so comparison on z\^{}r followed by
r-th roots makes them equal. Existence and completeness follow from the
finite affinoid field construction and the same finite-dimensional norm
argument.

The G action is transitive on these factors. Otherwise the sum of the
idempotents of one orbit would be a nontrivial invariant idempotent, in
contradiction with (C tensor E)\^{}G=C. Let H be the actual stabilizer
of y and let D\_y be its field factor. Invariant tuples in the
transitive product are equivalent to H-invariant elements of D\_y;
therefore

\[
 D_y^H=C.                                                   \tag{7}
\]

The map E-\textgreater D\_y is injective. The image of b has
\textbar H\textbar{} distinct H-conjugates, and D\_y=C{[}b{]}. The
polynomial product over those conjugates has coefficients in (7); the
minimal polynomial of b over C divides it and has every one of those
distinct conjugates as a root. Its degree is consequently exactly
\textbar H\textbar. This proves directly that D\_y/C is Galois with
group H. In particular the image of the finite separable field

\[
 F=E^H                                                     \tag{8}
\]

lies in C. If y is not fixed by all of G, H is proper and
\textbar H\textbar\textless\textbar G\textbar. This is a statement about
the actual point and its fibre, not a topological covering assertion.

\subsubsection{1.3. Purity, strictness, and ordinary
topology}\label{purity-strictness-and-ordinary-topology}

Strict affinoid presentations base change to strict ones. Inverse images
of ordinary affinoid interiors give goodness upstairs; separatedness is
preserved by base change of the closed diagonal, and boundarylessness
follows from foundation B. The same assertions hold for ordinary opens.

We check every minimal component, not just total dimension. In a strict
affinoid neighborhood M(A) of a pure n-dimensional space, each minimal
prime P has dim(A/P)=n: by spectral surjectivity a point of kernel P
exists, belongs to that component alone, and its local dimension is n.
Let J be the nilradical. It is nilpotent. The injection

\[
 A/J\longrightarrow\prod_{P\text{ minimal}}\operatorname{Frac}(A/P)
\]

remains injective on tensoring with E. Each field tensor is reduced by
separability, so (A/J) tensor E is reduced. The nilradical of A tensor E
is exactly J tensor E. The kernels of the finitely many maps of this
reduced tensor algebra to those field factors are primes whose
intersection is zero. A minimal prime Q contains their product and hence
one of them; minimality makes it that kernel. Its contraction is its
associated P. The finite integral injection

\[
 A/P\hookrightarrow (A\otimes_K E)/Q
\]

preserves dimension. For clarity, this last elementary fact follows from
lying over, going up, and incomparability: for a finite integral
injection of domains C0-\textgreater D0, localizing at p gives a nonzero
finite C0\_p-module. If pD0\_p=D0\_p, the determinant of I minus a
p-valued generator matrix is a unit annihilating it, a contradiction. A
maximal ideal over pD0\_p proves lying over. Apply the same argument
after quotienting by a prime to get going up. For incomparability,
quotient by the smaller prime and localize by the nonzero elements of
the base domain. The resulting finite integral domain over a field is a
field: a minimum-degree monic relation for a nonzero element expresses
its inverse using its nonzero constant term. Thus two strictly nested
primes cannot have the same contraction. Chain contraction and chain
lifting prove dimension equality.

Every minimal component upstairs therefore has dimension n, and the
local dimension formula proves purity. Reduction was used only in this
dimension calculation, not in any current, unit, or calibration step.

Inverse images of compact sets under pi are compact. Indeed a downstairs
compact set is covered by finitely many ordinary interiors of affinoid
neighborhoods. Its inverse image is closed in the finite union of their
compact scalar extensions. Hausdorffness ensures that the closedness
statements here are ordinary topological statements.

We will repeatedly use a full-fibre shrinking observation. Inside an
affinoid neighborhood V0 of x, let N be an ordinary open in V0\_E
containing the entire fibre over x. The compact complement has a closed
image in V0 not containing x. Choose a smaller strict affinoid
neighborhood V of x inside that image's complement and inside the
ordinary interior of V0. Then V\_E is contained in N. Its ordinary
interior gives the corresponding ordinary downstairs neighborhood. This
retains any finite collection of unit or strict norm conditions holding
on the whole fibre.

\subsection{2. A separable residue constant lifts to an actual
neighborhood}\label{a-separable-residue-constant-lifts-to-an-actual-neighborhood}

\textbf{Local lifting lemma.} Suppose F/K is finite separable and a
specified K-embedding F-\textgreater H(x) is given at a point of a good
strict K-space. There is a strict affinoid neighborhood V of x and a
bounded embedding F-\textgreater O(V) whose evaluation is the specified
embedding. Its restriction to an ordinary neighborhood U of x makes U an
F-analytic space on the same carrier and with the same actual analytic
function sheaf.

\textbf{Proof.} Let a be a primitive element of F, P its monic separable
minimal polynomial, and a also its specified image in H(x). Begin in a
strict affinoid neighborhood M(A0). The fraction field of A0/ker(x) is
dense in H(x). Thus choose g=f/b with b(x) nonzero and g(x) arbitrarily
close to a. Retain b invertible on a strict rational neighborhood, by an
inequality \textbar b\textbar\textgreater=\textbar k0\textbar{} strictly
satisfied at x, with k0 in K\^{}*. This turns g into an actual analytic
element. The following choices are made before the final shrink.

Fix R in \textbar K\^{}\emph{\textbar{} with
R\textgreater max(1,\textbar a\textbar), and fix d in
\textbar K\^{}}\textbar{} with
0\textless d\textless\textbar P'(a)\textbar. Put

\[
 M=\max\bigl(1,\max_{2\leq j\leq\deg P}
                     |p_j|R^{j-2}\bigr),
 \qquad P(T)=\sum_jp_jT^j.                                 \tag{9}
\]

Choose c in \textbar K\^{}*\textbar{} so small that

\[
 s:=c/d<1,\qquad q:=Mc/d^2<1,\qquad s< R.                   \tag{10}
\]

This is possible by taking powers of a K-element of norm \textless1.
Approximate a sufficiently closely that
\textbar g(x)\textbar\textless R,
\textbar P'(g)(x)\textbar\textgreater d,
\textbar P(g)(x)\textbar\textless c and also
\textbar g(x)-a\textbar\textless s. Preserve these strict inequalities
on the strict rational domain defined by

\[
 |g|\leq R,\quad |P(g)|\leq c,\quad |P'(g)|\geq d.            \tag{11}
\]

Here scalar representatives of R,c,d are chosen in K\^{}*. The rational
algebra is an actual complete quotient of a Tate algebra with variables
g/R, P(g)/c, and d/P'(g). The last relation makes P'(g) invertible as an
actual algebra element. Its quotient Banach norm has

\[
 \|g\|\leq R,\quad\|P(g)\|\leq c,
              \|P'(g)^{-1}\|\leq d^{-1}.                    \tag{12}
\]

These bounds arise from relations in the quotient algebra and remain
valid in the presence of nilpotents. We have not inferred a Banach bound
from the spectral seminorm.

For \textbar\textbar z\textbar\textbar\textless=R and
\textbar\textbar v\textbar\textbar\textless=min(1,R), binomial expansion
and the ultrametric inequality give

\[
 \|P(z+v)-P(z)-P'(z)v\|\leq M\|v\|^2,
 \qquad \|P'(z+v)-P'(z)\|\leq M\|v\|.                       \tag{13}
\]

Binomial coefficients have norm at most one; this also covers positive
characteristic. Starting with z0=g, iterate
z\_\{j+1\}=z\_j-P'(z\_j)\^{}\{-1\}P(z\_j). Let v\_j=z\_\{j+1\}-z\_j. The
first step has \textbar\textbar v0\textbar\textbar\textless=s. If
\textbar\textbar v\_j\textbar\textbar\textless=s, (13) and M
s/d=q\textless1 show that P'(z\_\{j+1\}) is invertible by a Neumann
series with inverse norm \textless=d\^{}-1. The Newton cancellation and
(13) give

\[
 \|v_{j+1}\|\leq (M/d)\|v_j\|^2\leq q\|v_j\|.              \tag{14}
\]

All z\_j have norm \textless=R. Thus the sequence converges in the
actual Banach algebra to z with P(z)=0,
\textbar\textbar z-g\textbar\textbar\textless=s, and invertible
derivative.

Evaluation is norm-contracting. Both z(x) and a are roots within s of
g(x). For two such roots r,r', the factorization of P(r)-P(r') gives

\[
 P(r)-P(r')=(r-r')\bigl(P'(g(x))+w\bigr),\quad |w|\leq Ms<d.
\]

The factor in parentheses is nonzero; hence r=r'. Consequently z(x)=a.
The homomorphism K{[}T{]}/P-\textgreater O(V) sending T to z is unital
and injective, because its source is a field and its target is nonzero
(it has x). It is bounded: finite-dimensional norm equivalence bounds
the coefficient maps in the basis 1,a,\ldots,a\^{}\{deg P-1\}, and their
images are bounded elements of O(V). This proves the desired lift.

We verify the F-analytic interpretation rather than treating it as a
change of notation. Let A=O(V). A bounded embedding of the finite valued
field F gives
\textbar\textbar u\textbar\textbar\_A\textless=C\textbar u\textbar{} for
u in F. The equivalent complete norm

\[
 \|b\|_F:=\sup_{u\in F^*}\|ub\|_A/|u|
\]

is F-homogeneous and submultiplicative: the original norm is bounded
above by this norm, and
\textbar\textbar ab\textbar\textbar\_F\textless=\textbar\textbar a\textbar\textbar\_F
\textbar\textbar b\textbar\textbar\_A\textless=\textbar\textbar a\textbar\textbar\_F\textbar\textbar b\textbar\textbar\_F.
Choose a strict K-Tate presentation of A. Its finitely many generators
are power-bounded in the original quotient norm and hence in the
equivalent F-norm. They define a bounded surjection
F\{T1,\ldots,Tq\}-\textgreater A. The quotient norm is equivalent to the
F-norm: boundedness gives one inequality; lifting through the original
K-Tate presentation gives the other, since K\{T\} is a subalgebra of
F\{T\}. Thus A is strictly F-affinoid.

Every bounded K-multiplicative seminorm on A restricts to the unique
extending absolute value on F, so the K- and F-affinoid spectra coincide
with their topology. Rational domains use the same elements of A and the
same inequalities; their actual quotient algebras are the same with the
induced bounded F-action. These common rational bases identify the
analytic sheaves on ordinary opens. Restrict to an ordinary neighborhood
U of x inside V to obtain the assertion. Nilpotent elements have been
retained throughout. \textbf{End of proof.}

The same argument identifies ordinary form sheaves on U over K and F.
Their unit germs, compact moment images, and ordinary descriptions are
identical. Foundation C's compact-image nullity criterion identifies
their relations. Supports, wedge products, d' and d'\,' therefore agree.
Equation (3) and the definitions of top measures give, on this same
carrier and for every complementary test alpha,

\[
 \int_U h\,(D\alpha)_{\mathrm{Rad},K}
 =[F:K]\int_U h\,(D\alpha)_{\mathrm{Rad},F}.                 \tag{15}
\]

The factor is a positive real integer, even in wild characteristic, and
can be cancelled for vanishing.

If U was chosen ordinary open in X, it is boundaryless over K. The
composition equality for U-\textgreater M(F)-\textgreater M(K), whose
second map is finite separated, makes U boundaryless over F as well. The
same affinoid algebras and minimal primes prove purity over F. Goodness
and strictness were checked above. Separatedness is retained: the
F-fibre product is the closed equalizer, inside the K-fibre product, of
the two F-actions, and the K-diagonal already is a closed immersion.
Thus U is an admissible input to the induction over F.

\subsection{3. The off-fixed-point
induction}\label{the-off-fixed-point-induction}

We prove the forward assertion of (T) by strong induction on
\textbar G\textbar, simultaneously for every field and every space in
its stated scope. The trivial group is the identity assertion. Suppose
smaller groups have been treated. Let

\[
 S=\{y\in X_E:g y=y\text{ for all }g\in G\}.
\]

It is closed, as a finite intersection of equalizers of continuous maps
on a Hausdorff space. Let y not belong to S and put x=pi(y). Its group H
is proper. The specified F=E\^{}H embedding in H(x) from section 1.2
lifts on U by section 2. By (15), the restriction of h has zero
F-Hessian current on U.

For its actual F-algebra A, separability and normality give

\[
 F\otimes_K E\simeq\prod_{\tau:F\hookrightarrow E}E,
 \qquad
 A\otimes_K E\simeq\prod_{\tau}A\otimes_{F,\tau}E.           \tag{16}
\]

To see the first identity directly, a primitive polynomial for F/K
splits in E into its distinct linear factors; use the Chinese remainder
theorem. The second follows by tensoring over F. These are identities of
the actual algebras, not their reductions. Their finite-module
topologies agree by section 1.1, and their idempotents give
open-and-closed analytic components. The component selected by y is the
one for the specified embedding. On each component, regard U as a
tau(F)-space via tau. E/tau(F) is Galois with a group conjugate to H, of
order \textbar H\textbar\textless\textbar G\textbar. All the hypotheses
of the induction have been checked in sections 1-2. It proves zero
current on that component, in particular near y.

Consequently B\_\{h\_E\} vanishes on every complementary test supported
in X\_E minus S. For a compact such support, finitely many ordinary
neighborhoods of the kind just constructed suffice; compact smooth
partitions (proved next) split the test among them. There is no current
pullback in this step.

\subsection{4. Compact partitions and continuous scalar
approximation}\label{compact-partitions-and-continuous-scalar-approximation}

One can reconstruct the compact-cutoff input directly here, avoiding any
overbroad partition statement. Given a point x inside an ordinary open
U, take an affinoid neighborhood V of x and put F=V minus (its ordinary
interior intersect U). This is compact and does not contain x. Analytic
functions on V separate x from every point q in F by their norms. For
such a function f, compose log\textbar f\textbar{} on its nonzero locus
with a real smooth function that takes the desired distinct constant
values near \textbar f(x)\textbar{} and \textbar f(q)\textbar. If either
norm is zero, choose the composing function constant on a sufficiently
negative logarithmic half-line; extension by that constant across the
zero locus is ordinary smooth. Thus there is an ordinary smooth function
on V equal to one near x and zero near q. A finite subcover of F and the
product of these functions gives b on V, one near x and zero on a
relative neighborhood of F. Extend it by zero outside V. It is smooth on
the ordinary interior of V and outside V; at each point of F, a global
ordinary neighborhood meets V only where b is zero, so it is smooth
there too. The support is a closed subset of compact V contained in its
interior intersect U. For a compact C in U, finitely many such bumps and
a smooth real composition of their sum of squares give a compact smooth
cutoff equal to one near C. This proves exactly the compact cutoff used
below. The embedded primary statement 15.2.5 is consistent with this
proof; the printed sign typo in the one-point construction of A.2.3 is
not an additional premise.

For a compact C and a finite ordinary open cover U1,\ldots,Uk of C, use
foundation D point by point to obtain smooth compact bumps equal to one
near their centers and supported in the appropriate Ui. Compose each
with a nonnegative smooth real function equal to zero near zero and one
near one. Finitely many bumps cover C. Group them by index into psi\_i
and set s=sum psi\_i. It is positive near C. A second compact bump chi
supported in \{s\textgreater0\}, equal to one near C, gives

\[
 \theta_i=\chi\psi_i/s\quad\text{on }\{s>0\},
 \qquad \theta_i=0\quad\text{elsewhere}.                    \tag{17}
\]

These are ordinary smooth, compactly supported inside Ui and sum to one
near C. Extension by zero is smooth because the support lies compactly
inside the defining ordinary open. The same observation applies to every
test extension by zero below. No global paracompactness is required.

Given any continuous h on a fixed compact C and epsilon\textgreater0,
take a finite ordinary neighborhood cover on which the oscillation of h
is \textless epsilon. Apply (17), choose a value c\_i of h on each
member meeting C, and put a=sum c\_i theta\_i. Then a is an ordinary
smooth compact scalar function and \textbar a-h\textbar\textless epsilon
on C. Thus the exact scalar approximation needed below follows from
compact cutoffs; it is not a tropical factorization of h.

\subsection{5. Norm forms match every test on the fixed moment
subspace}\label{norm-forms-match-every-test-on-the-fixed-moment-subspace}

Let beta be an arbitrary compact complementary ordinary test on X\_E.
For y in S intersect supp(beta), the whole fibre over x=pi(y) is the
singleton \{y\}: transitivity in section 1.2 proves this. An ordinary
description of beta near y uses actual units f1,\ldots,fk and an ambient
smooth superform. Intersect its finitely many G-translates so every g
f\_i is an actual unit there. Full-fibre shrinking gives a strict
affinoid neighborhood V of x whose whole scalar extension is inside that
ordinary unit domain.

On V\_E take the orbit moment

\[
 t=(t_{g f_i})_{i,g}: V_E\longrightarrow\mathbb{R}^{k|G|}.
\]

Repeated units may be retained as separate coordinates. G permutes these
coordinates. Let W be its fixed vector subspace, let
i:W-\textgreater R\^{}\{k\textbar G\textbar\} be inclusion, and let
P0=\textbar G\textbar\^{}-1 sum\_g g be the real linear averaging
projection. At every actual G-fixed analytic point t belongs to W. The
coordinate corresponding to each orbit in P0 t is

\[
 |G|^{-1}\sum_g t_{g f_i}=|G|^{-1}t_{N(f_i)}\circ\pi.       \tag{18}
\]

By (6) these norms are actual downstairs units on V. The denominator in
(18) is a real denominator in a superform, not division in K.

Let omega be the ambient coefficient superform describing beta in this
orbit moment (pull back its original coefficient form under coordinate
projection). At t(y) the point is fixed. Shrink V again so that both
t(V\_E) and iP0t(V\_E) lie in the ambient coefficient domain. By a
Euclidean compact bump, extend omega to a global smooth ambient form,
unchanged on a neighborhood of both compact sets. This elementary
extension uses a compact subset of an open Euclidean set and introduces
no singularities. The form

\[
 \omega_0=P_0^*i^*\omega
\]

has i\^{}\emph{(omega-omega0)=0. Equation (18) makes t\^{}}omega0
exactly the pullback of a downstairs ordinary smooth form alpha0. This
is matching on W only; it does not say beta, or even an invariant beta,
descends.

The compact set T=pi(S intersect supp(beta)) is covered by finitely many
ordinary downstairs neighborhoods Vi of this kind. Choose compact
downstairs theta\_i by (17) with sum theta\_i=1 near T. Then

\[
 \beta=\beta_{\rm rem}+\sum_i\beta_i,
 \quad \beta_i=\pi^*\theta_i\,\beta,
 \quad \beta_{\rm rem}=(1-\pi^*\sum_i\theta_i)\beta.          \tag{19}
\]

The remainder has compact support disjoint from S, so section 3 kills
its pairing. For each i define

\[
 \beta_{0,i}=\pi^*(\theta_i\alpha_{0,i}),
 \qquad \gamma_i=\beta_i-\beta_{0,i}.                        \tag{20}
\]

All these forms are extended by zero from their ordinary chart. They
have compact support contained in one fixed compact pi\^{}-1(supp
theta\_i) inside that chart. Foundation F, applied to the differentiated
downstairs test theta\_i alpha0\_i, gives

\[
 B_{h_E}(\beta_{0,i})=B_h(\theta_i\alpha_{0,i})=0.            \tag{21}
\]

No degree factor and no unexplained continuous-current pullback occurs.
It remains to kill gamma\_i. Its coefficient-form restriction to W is
zero; if theta\_i needs extra local moment coordinates, use W times the
entire space of those extra coordinates. Multiplication by theta\_i
preserves this restriction-zero assertion.

\subsection{6. The full tensor collar estimate, including boundary-face
contact}\label{the-full-tensor-collar-estimate-including-boundary-face-contact}

Here is the elementary estimate needed for (20).

\textbf{Collar lemma.} Let Q be a fixed linear map on a Euclidean space
V with kernel W; an orthogonal projection may be used. Let sigma be a
compact n-dimensional affine polytope. Let xi be a smooth (n-1,n-1)
superform on a neighborhood of sigma and of the nearby W points, with
i\^{}*xi=0 on W. Take a fixed smooth radial function chi on Q(V), equal
to one near zero and supported in the unit ball, and define

\[
 \xi_\delta(z)=\chi(Qz/\delta)\xi(z),\quad 0<\delta\leq1.
\]

On sigma, the coefficient L1 norms of D xi\_delta are uniformly bounded
as delta tends to zero. The coefficient L1 norm of xi\_delta tends to
zero. These conclusions still hold with arbitrary fixed constant cell
volume and with additional Euclidean coordinates on which Q is zero.
Only the restriction of the superforms to the affine tangent of sigma is
measured.

\textbf{Proof.} Let L be the n-dimensional tangent plane of sigma and
put r=rank(Q\textbar L). If its affine span misses W, the affine image
Q(span sigma) is a closed affine subspace not containing zero, at
positive distance from it. For sufficiently small delta the collar is
empty. Otherwise choose w0 in that affine span intersect W and write
z=w0+u+v, with u in L intersect W and v in a fixed complement L1 of
dimension r. Q\textbar L1 is an isomorphism to its image. Compactness
bounds u in one fixed box, and \textbar Qz\textbar\textless=delta bounds
v in a box of side O(delta). Fubini and the fixed determinant give
collar volume O(delta\^{}r). This argument does not require w0+u to
belong to sigma: it works for a cell meeting W only at a boundary face,
or with zero-dimensional intersection. If a compact cell itself misses
W, its distance is positive and the same conclusion follows.

For r=0 a cell intersecting W lies in W. Its xi restriction is
identically zero, so the restriction of D xi\_delta is also identically
zero. For r\textgreater=2, cutoff derivatives of orders j=0,1,2 are
O(delta\^{}-j) on the fixed cell, and xi and its first two coefficient
derivatives are bounded. The Leibniz expansion therefore bounds D
xi\_delta by O(delta\^{}-2). Multiplying by O(delta\^{}r) gives a
uniform integral bound.

For r=1 write the affine tangent coordinates as
(u1,\ldots,u\_\{n-1\},t), with z=w0+u+t v and v fixed, Qv nonzero. The
cutoff restricted to this plane depends only on t and has derivatives
O(delta\^{}-1) and O(delta\^{}-2). Expand the restricted xi in the
exterior basis using d'u\_j,d't and d'`u\_j,d'`t. There is exactly one
possible basis monomial of bidegree (n-1,n-1) containing neither d't nor
d''t:

\[
 A(u,t)\,d'u_1\wedge\cdots\wedge d'u_{n-1}
          \wedge d''u_1\wedge\cdots\wedge d''u_{n-1}.         \tag{22}
\]

Its coefficient A(u,0) is zero, because the u-plane at t=0 is tangent to
W and i\^{}*xi=0 there. Smoothness on a fixed compact neighborhood gives
\textbar A(u,t)\textbar\textless=C\textbar t\textbar. The comparison
point w0+u can lie outside the cell; it is still on W and within
O(delta) of its compact collar. This is exactly why restriction-zero is
needed on the ambient W, not just at points of sigma intersect W.

The second-cutoff term is chi'`(t) d't wedge d''t wedge xi. Every
monomial of xi containing at least one transverse generator is
annihilated by a repeated exterior generator. Only (22) remains, bounded
by O(delta\^{}-2) O(delta)=O(delta\^{}-1) on the collar. All
first-cutoff terms are O(delta\^{}-1), since their other coefficients
and derivatives are bounded. The term chi D xi is bounded. The collar
volume is O(delta), so all differentiated terms have bounded coefficient
L1 norm. This includes n=1: then (22) is just the scalar coefficient,
vanishing on W.

For the undifferentiated form, coefficients are bounded and every
positive-rank collar has volume tending to zero; rank-zero cells either
miss the collar or give the zero restriction. A finite number of cells,
fixed coordinate Jacobians, and fixed positive calibration constants do
not change either conclusion. \textbf{End of proof.}

This lemma is not a claim that a small collar gives small differentiated
variation; the latter is only bounded and need not tend to zero. Rank
one uses both exterior-tensor cancellation and the first-order vanishing
in (22). No balanced-fan recognition theorem or curve probe is involved.

\subsubsection{6.1. From coefficient bounds to actual top-measure
variation}\label{from-coefficient-bounds-to-actual-top-measure-variation}

On a compact affinoid moment chart of dimension n, foundation E gives a
finite polyhedral image and a finite skeletal n-cell refinement. Above
each image cell sigma, the finitely many skeletal cells map by affine
isomorphisms and have positive constant calibration d\_C times a
primitive volume. For an ambient top form zeta, its density on each such
sheet is the same affine-tangent scalar density, multiplied by that
positive constant. Thus taking absolute values before summing gives

\[
 \|(t^*\zeta)_{\mathrm{Rad}}\|_{\mathrm{TV}}
 \leq\sum_{\sigma} c_\sigma
        \int_\sigma\|\zeta|_{T\sigma}\|\,d\operatorname{vol}_\sigma,
                                                                  \tag{23}
\]

for fixed finite c\_sigma. In an oriented tangent basis the
scalar-density version is an equality; the norm inequality allows
arbitrary bases. There is no signed cancellation hidden in this bound.
It is the local construction in CLD 5.8.4-5 plus 14.3.2 and 14.5.2, not
an inference from signed pushforward equality alone. The constants
include all actual finite-flat fibre dimensions, and therefore all
nilpotent lengths. Cells below dimension n contribute zero. If the whole
moment image has dimension \textless n, every represented top form is
zero. We never replace that case by an assumed universal dimension
equality.

Finitely many such compact charts covering a fixed ordinary compact
support suffice. Top measures restrict compatibly to domains;
overcounting their positive variations in an overlapping cover yields an
upper bound. No derivative of a partition is needed to obtain (23).

\subsection{7. Weak convergence on the full continuous
multiplier}\label{weak-convergence-on-the-full-continuous-multiplier}

Fix a gamma=gamma\_i from (20), its orbit moment t, and put Q=1-P0.
Choose the radial chi of section 6 and define, in its chart and extended
by zero,

\[
 e_\delta=\chi(Qt/\delta)\gamma.                            \tag{24}
\]

These are ordinary smooth complementary tests with supports in the same
compact C. The compact test (1-chi(Qt/delta)) gamma has support disjoint
from S: on S, Qt=0 and the multiplier equals zero on an ordinary
neighborhood. Section 3 therefore gives the \textbf{exact} identity

\[
 B_{h_E}(\gamma)=B_{h_E}(e_\delta)\quad(\delta>0).             \tag{25}
\]

Cover C by finitely many strict affinoid neighborhoods contained in the
chart, and add moment coordinates for local descriptions of the fixed
downstairs cutoff theta\_i. On each enlarged ambient space, gamma is a
smooth coefficient form whose restriction to W times the extra space
vanishes. The collar lemma, (23), and the finite cover give

\[
 \sup_{0<\delta\leq\delta_0}
       \|(D e_\delta)_{\mathrm{Rad},E}\|_{\mathrm{TV}}<\infty. \tag{26}
\]

The constant may depend on this test and this compact carrier. It is
independent of delta. This is a statement about smooth differentiated
forms, not a Radon regularity hypothesis on D{[}h\_E{]}.

Let a be any fixed ordinary smooth compact scalar multiplier. Stokes
twice, on the boundaryless pure space X\_E, gives

\[
 \int_{X_E}a\,D e_\delta=\int_{X_E}e_\delta\wedge D a.        \tag{27}
\]

For explicit signs, 0=integral d'(a d'`e) says integral aDe= -integral
d'a wedge d'`e. Next 0=integral d'`(d'a wedge e) says integral d'a wedge
d'`e=integral d''d'a wedge e=-integral Da wedge e. The total degree of e
is even, so Da wedge e=e wedge Da. These calculations require no
symmetry of e and thus treat every ordinary complementary test.

To show the right side tends to zero, add \textbf{the moment coordinates
describing a} on a finite compact chart cover of C. This additional
refinement is essential: L1 information on a smaller moment need not by
itself control a wedge with Da. The enlarged fixed subspace is W times
all the added coordinates, and the coefficient restriction of gamma
still vanishes there. The undifferentiated conclusion of the collar
lemma gives L1 convergence to zero on every calibrated n-cell of this
enlarged moment; Da has bounded coefficients on the fixed compact
images. Formula (23) applied to e\_delta wedge Da proves

\[
 \int a\,(D e_\delta)_{\mathrm{Rad},E}\longrightarrow0.       \tag{28}
\]

The refinement and constants for (28) may depend on a. They do not need
to be uniform in a. Uniformity in delta for (26), already obtained
without a, is what the continuous passage uses.

Approximate h\_E on C within epsilon by an ordinary smooth compact a
using section 4. If M is the bound in (26), then

\[
 \limsup_{\delta\to0}
 \left|\int h_E\,(D e_\delta)_{\mathrm{Rad},E}\right|
 \leq M\varepsilon.                                       \tag{29}
\]

Let epsilon tend to zero. Equations (25) and (29) imply
B\_\{h\_E\}(gamma)=0. By (19)-(21), B\_\{h\_E\}(beta)=0 for the original
arbitrary compact test. This completes the induction and the forward
implication of (T).

\subsection{8. Reverse transport and the normalization with actual
lengths}\label{reverse-transport-and-the-normalization-with-actual-lengths}

For alpha in A\_c\^{}\{n-1,n-1\}(X), pi\^{}*alpha is an ordinary smooth
complementary test with compact support, since pi has compact inverse
images. It is an allowed test upstairs, not a restriction of the
upstairs test class to pullbacks. Apply (2) to omega=Dalpha and use the
continuous multiplier h and the defining pushforward of a signed Radon
measure:

\[
 \begin{split}
 B_{h_E}(\pi^*\alpha)
 &=\int_{X_E}(h\circ\pi)\,(\pi^*D\alpha)_{\mathrm{Rad},E}\\
 &=\int_X h\,(D\alpha)_{\mathrm{Rad},K}=B_h(\alpha).
 \end{split}                                               \tag{30}
\]

This proves the reverse implication with no {[}E:K{]} factor.

Here is the finite-length algebra in the imported scalar calibration
proof. If R is a finite Artinian C-algebra of dimension d and C' is a
complete extension, then R tensor\_C C' is still d-dimensional as a
C'-vector space. Its decomposition into local Artinian factors R\_j has

\[
 d=\sum_j\dim_{C'}R_j
   =\sum_j\operatorname{length}_{R_j}(R_j)
                         [\kappa_j:C'].                    \tag{31}
\]

To verify the second equality, take a composition series of each local
factor; every simple quotient is its residue field. Tensoring a C-basis
proves the first dimension statement, without any reducedness or
separability assumption on R. A finite-flat moment fibre is exactly such
an actual algebra. Above the corresponding standard torus Gauss point,
scalar extension chooses its canonical Gauss lift and partitions those
fibre factors; (31) says their degrees sum to the original degree. On
every skeletal n-cell this is precisely the primary proof of CLD 14.3.6,
using CLD 13.3.3, and gives (2). The source appendix includes those
complete proofs. In contrast, restricting an F-moment to K composes its
finite-flat torus projection with
G\_m,F\textsuperscript{n-\textgreater G\_m,K}n, whose actual fibre rank
is {[}F:K{]}; this gives (3), not (2).

At no point did we replace a local Artinian algebra by its residue field
or a singular component by a component of maximal total dimension.
Equations (4), (5), (6), (11)-(16), (23), and (31) all preserve the
actual structures required in the assignment.

\subsection{Audit boundaries and comparison with the
catalogue}\label{audit-boundaries-and-comparison-with-the-catalogue}

The new scope is every finite Galois scalar extension on every pure
n-dimensional space in the original strict, good, separated,
boundaryless setting, for finite continuous multipliers. It includes
non-quasi-smooth germs in all positive dimensions and actual nilpotent
thickenings. Unlike the catalogue's scalar-stable special cases, its
upstairs current equation is a conclusion rather than a hypothesis. It
is also stronger than a curve-only scalar-test result: the proof starts
with an arbitrary compact (n-1,n-1) form, uses full tensor coefficients,
and treats every local n-cell and lower-dimensional moment image.

The conclusion is still just a missing gate. The original recognition
target also requires a bridge from continuous current positivity of the
two signs to the classical curve-test notion (or another recognition
argument). No such bridge is proved. No completed-algebraic ascent is
claimed: finite-stage descent and gluing of arbitrary complementary
ordinary forms would have to be established separately. No mixed-finite
or transcendental ascent is claimed. The unrestricted non-strict
arbitrary-radius statement from the catalogue is untouched.

The complete authoritative argument, both canonical native notes,
all thirty-three embedded primary units, author audit and research limits
were read before publication. The independent Xean native Codex review
used \texttt{gpt-6.1-sol} with \texttt{xhigh} reasoning and returned
PASS for this precise partial theorem. Its full report, all twenty-seven
premises and all forty passages were read; passage indices and exact
statement matches were checked. Frozen proof, source, status and
primary-body bindings accompany the review.

The primary-source quotations, with their complete supplied proofs,
remain in the accompanying frozen argument. The following harmless
source-display errata were already avoided by the mathematical proof:
graded commutativity has exponent $(p+q)(p'+q')$; the one-point cutoff
is one near the first point and zero near the second; the comparison in
CLD 5.10.1(a) is equality; the finite case in 13.3.3 has
$\dim_y(Y)=0$; pushed volume in 15.3.1 equals downstairs volume; and
11.4.1 uses the inverse image of $\operatorname{Int}(Y/X)$.
Raw frozen quotations are preserved unchanged. Relative-interior base
change is used only as preservation of interior points, as stated above.

The standalone source was opened in the native LaTeX editor for
compilation. Same-model review, finite diagnostics, source integrity and
compilation are not formal verification or priority certification.
The full original recognition problem remains open.

\subsection{Sources and accompanying material}

The accompanying complete frozen argument embeds the essential primary
units from Chambert-Loir and Ducros, \emph{Formes diff\'erentielles
r\'eelles et courants sur les espaces de Berkovich},
\url{https://arxiv.org/abs/1204.6277v2}; Ducros,
\emph{Espaces de Berkovich, polytopes, squelettes et th\'eorie des
mod\`eles}, \url{https://arxiv.org/abs/1203.6498v3}; and Berkovich,
\emph{\'Etale cohomology for non-Archimedean analytic spaces},
Publ. Math. IHES 78 (1993). The independently checked CLD composition
statement supplies the boundary facts needed here. The exact imported
foundations A--G and every load-bearing new proof above remain in this
typeset paper. Primary quotations are retained in the frozen argument
rather than duplicated in this rendering.

\end{document}
